\section{引言}
开放词汇语义分割通过文本指定目标类别，为遥感图像中多样化地物的识别提供了灵活接口。与自然图像相比，遥感图像中的尺度差异、密集分布及相似地表纹理增加了像素级分类的难度。如何结合遥感语义先验与精细空间表征，是将视觉语言模型用于该任务的重要问题。

SAM 3 能够根据文本概念生成分割响应，为开放词汇遥感解译提供基础~\cite{carionSAM3Segment2026}。在逐类构建语义预测的流程中，各类别条件特征独立形成，解码前的类别关系缺少显式建模。道路与裸地等类别的局部外观可能相似，仅依靠单一类别提示难以充分利用同一场景中的语义上下文。

本文提出 ReCA-SAM3，将 RemoteCLIP 的遥感图文先验引入 SAM 3 的解码前特征聚合。该方法以局部图文相似度初始化语义状态，通过类间注意力交换不同提示的信息，并以类内窗口注意力建模空间联系。语义—细节融合模块进一步结合聚合特征、原始类别条件特征与多尺度空间表征，复用冻结的 SAM 3 解码路径。本文围绕跨类别交互、遥感语义引导及方向融合设置模块对照，并在八个遥感数据集上组织源域与跨数据集评测。

